\subsection{Robustness}\label{app:robustness}

This subsection states and proves Proposition~\ref{prop:control} and the results that Section~\ref{sec:robustness}
summarizes. Each proof uses the conventions of Appendix~\ref{app:beliefs}.

\subsubsection*{The fixed-information control}

\begin{proposition}[Fixed information]\label{prop:control}
Fix the joint law of the class \(j\) and a signal \(S\) that the challenger observes.
Fix the class value laws \(G_j\) and the cost law. In each economy \(i\), require that, conditional on \((j,S)\), the joint law of \((\theta,R)\) is \(G_j\otimes F_i\). The cost is independent of the class, both values and \(S\). If
the incumbent strengthens from \(F_0\) to \(F_1\), entry is nonincreasing when an indifferent high-cost challenger prepares. Its decrease is exactly
\[
\Ent_0-\Ent_1=\Pr\bigl\{B_{F_1}(\mu(S))<C\le B_{F_0}(\mu(S))\bigr\},\qquad \mu(S)=\Pr(j=H\mid S),
\]
and the decrease is strict exactly when this probability is positive.
\end{proposition}

\begin{proof}[Proof of Proposition~\ref{prop:control}]
The cost is independent of the class, the values and the signal \(S\), so \(\mu(S)=\Pr(j=H\mid S)\) is the
challenger's complete class belief at its decision. The conditional joint law \(G_j\otimes F_i\) gives gross profit \(B_{F_i}(\mu(S))\) in economy \(i\).
When an indifferent high-cost challenger prepares, the challenger prepares exactly when
\(C\le B_{F_i}(\mu(S))\). Proposition~\ref{prop:returns}(a) gives \(B_{F_1}(\mu)\le B_{F_0}(\mu)\) at every belief.
Hence, pointwise,
\[
\1\{C\le B_{F_0}(\mu(S))\}-\1\{C\le B_{F_1}(\mu(S))\}=\1\{B_{F_1}(\mu(S))<C\le B_{F_0}(\mu(S))\}\ge0 .
\]
The joint law of \((S,C)\) is the same in both economies. Take expectations: \(\Ent_0-\Ent_1\) equals the
probability in the proposition, which is nonnegative, and it is positive exactly when the decrease is strict. A
tie probability \(q\) that is the same in both economies assigns the preparation probability
\(\1\{C<B_{F_i}\}+q\,\1\{C=B_{F_i}\}\). With \(B_{F_1}\le B_{F_0}\), the difference of the two expressions is \(0\) for
\(C<B_{F_1}\) and for \(C>B_{F_0}\), \(1-q\) at \(C=B_{F_1}<B_{F_0}\), \(1\) strictly between, \(q\) at \(C=B_{F_0}>B_{F_1}\),
and \(0\) when the two profits are equal; it is nonnegative, so entry is nonincreasing under that common tie probability.
\end{proof}

\paragraph{The control in the slot.}
The experiment fixes \(\lambda\), \(\beta\), \(\pi\), the cost law and the full orders
\((\sigma_H,\sigma_L)=(\delta_1,\delta_{-1})\). Lemma~\ref{lem:beliefs}(d) gives the posteriors \((m,\pi,M)\). Let
\(e_r(\mu)=\rho+(1-\rho)\1\{B_r(\mu)\ge c_H\}\) and \(P(x)=t_0+e_r(\mu(x))\,w(\mu(x))\). This entry is nondecreasing
in the posterior, so by the paragraph on constructed tuples in Appendix~\ref{app:beliefs} the three prices differ,
and each price reveals its posterior. The challenger's signal is therefore the order, with a law that does not
depend on \(r\). Given the class and order, the value law is \(G_j\otimes F_r\), as Proposition~\ref{prop:control} requires.
In the example, \(\tau(r)=[2r c_H-(\ell^2-p^2)]/(2rh-r^2-\ell^2)\) rises and \(\tau(\ell)=c_H/(h-\ell)>\pi\), so only the
top posterior \(M\) can recruit, and it recruits exactly when \(r\le r_C\). Entry is therefore
\(\rho+(1-\rho)(b+\lambda/2)=\reg{p4_frozen_entry_below_ceiling}\) on \((\ell,r_C]\) and
\(\rho=\reg{p4_frozen_entry_above_ceiling}\) above \(r_C=\reg{rC_B}\). The probability of
Proposition~\ref{prop:control} is \(\reg{p4_control_crossing_r0_r1}\) from \(r_0\) to \(r_1\), because the top posterior
recruits at both strengths, and \(\reg{p4_control_crossing_r1_r2}\) from \(r_1\) to \(r_2\). At \(r_0\) the frozen profile
is a control and not an equilibrium: its purchase and sale earn \(\Delta_0(1-M)-k<0\) and \(\rho\Delta_0m-k<0\). At
\(r_1\), the bound \(\rho\kappa\Delta_1>k\) of Step 2 of the proof of Proposition~\ref{prop:ccc} forces full orders
in every candidate, so the frozen tuple is the unique equilibrium there.

When the challenger does not see the price, its belief stays at \(\pi\), and \(B_{r_i}(\pi)<c_H\) gives entry
\(\rho\) after every order at both strengths. At \(r_0\), \(\Delta_0<k\) forces zero orders; at \(r_1\), the bound
\(\rho\kappa\Delta_1>k\) holds with \(e\equiv\rho\) and forces full orders. The prices are
\(t_0+\rho\,w(\mu(x))\), and every condition holds by the computations of Steps 1 and 3 of the proof of
Proposition~\ref{prop:ccc} with \(e\equiv\rho\). Entry is \(\reg{p4_hidden_r0_entry}\) at both strengths.
Table~\ref{tab:outcomes} lists all four outcomes.

\subsubsection*{Preparation cost laws}

For this robustness result, let the preparation cost \(C\) be independent of
the class, the value, the incumbent's value and the slot, with the law
\[
\Phi=\rho \Phi^0+(1-\rho)\Phi^1,\qquad \operatorname{supp}\Phi^0\subset[0,\bar c_0],\qquad \operatorname{supp}\Phi^1\subset[c_-,c_+].
\]
Each component can be atomless or have atoms. A challenger prepares when gross profit at its belief exceeds its
cost, and it stays out when gross profit falls short. The low component plays the role of the zero cost, and the high
component plays the role of \(c_H\).

With a general cost law, the preparation rule assigns \(\zeta(\bar P,C)\in[0,1]\) to each price and cost, and (E3)
reads: \(\zeta(\bar P,C)=1\) if \(B(\bar\mu(\bar P))>C\), and \(\zeta(\bar P,C)=0\) if \(B(\bar\mu(\bar P))<C\). Entry
after the order \(x\) is \(e(x)=\int\zeta(P(x),C)\,d\Phi(C)\). Every step of Appendix~\ref{app:beliefs} that does not
use the two-point law carries over; in particular, Lemma~\ref{lem:trading} holds for any entry in \([\rho,1]\).

\begin{proposition}[Preparation cost laws]\label{prop:atomless}
Let the classes satisfy the class ranking~\eqref{eq:classes}, and let \(F_1\) strengthen \(F_0\). Suppose
\[
0<\bar c_0<B_1(m),\qquad B_0(\pi)<c_-<c_+<B_1(M),\qquad \Delta_0<k<\rho\kappa\Delta_1.
\]
\begin{enumerate}[label=(\roman*),leftmargin=2em]
\item At \(F_0\), the unique orders are zero, and entry is \(\rho\).
\item At \(F_1\), the unique orders are full, and only a purchase recruits the high component. Entry is
\(\rho+(1-\rho)(b+\pi\lambda)\). High-class ownership rises if and only if
\([\rho+(1-\rho)(b+\lambda)]\,W_H(F_1)>\rho\,W_H(F_0)\).
\end{enumerate}
\end{proposition}

\begin{proof}[Proof of Proposition~\ref{prop:atomless}]
\emph{Step 0: consequences of the hypotheses.} The class ranking~\eqref{eq:classes} gives \(\Delta_i\ge0\) and
\(w_L\ge0\). Since \(k>0\), \(k<\rho\kappa\Delta_1\) gives \(\Delta_1>0\). The support gaps and
Proposition~\ref{prop:returns}(a) give \(B_1(M)>c_+>c_->B_0(\pi)\ge B_1(\pi)\). Since \(M>\pi\),
\(B_1(M)-B_1(\pi)=(M-\pi)(g_{H,1}-g_{L,1})>0\), so \(g_{H,1}>g_{L,1}\). The proof needs no condition \(\Delta_0>0\).

\emph{Step 1: the floor at the nonzero orders.} Fix a candidate in either economy. A price reached by a
nonzero order has its posterior in \([m,M]\), so
\[
B_{F_i}(\bar\mu)\ge B_{F_i}(m)\ge B_1(m)>\bar c_0.
\]
Every cost in the support of \(\Phi^0\) therefore prepares strictly at both nonzero-order prices, giving entry
at least \(\rho\) after a purchase and after a sale.

\emph{Step 2: orders and preparation by the low component.} At \(F_0\), entry in \([0,1]\) gives correct gross
returns at most \(\Delta_0<k\), and wrong-signed orders earn at most \(-k\). Both classes therefore abstain.
Every order and price has posterior \(\pi\), and
\(\bar c_0<B_1(m)\le B_0(\pi)<c_-\) gives entry \(\rho\).
At \(F_1\), the nonzero-order floor and the trading bound give correct gross returns at least
\(\rho\kappa\Delta_1>k\), so both classes place full orders. Their posteriors are \((m,\pi,M)\), and each price
belief is an average of these. The low component therefore prepares strictly at every price. These order
comparisons hold under both investor tie rules. The price formula and the strict support windows give the
equilibrium tuples and outcomes of Proposition~\ref{prop:atomless}.


\emph{Step 3: the tuples.} At \(F_0\): zero orders, \(\mu\equiv\pi\), \(\bar\mu\equiv\pi\), the single price
\(t_{0,0}+\rho\,w_0(\pi)\), and \(\zeta(\bar P,C)=\1\{C\le B_0(\pi)\}\). At \(F_1\): full orders, posteriors
\((m,\pi,M)\), \(\zeta(P(x),C)=\1\{C\le B_1(\mu(x))\}\), entry \((\rho,\rho,1)\), prices
\(t_{0,1}+e(x)\,w_1(\mu(x))\) and \(\bar\mu(P(x))=\mu(x)\); truthful bids in both. Equilibrium conditions: (E1) and (E2) hold as
in Steps 1 and 3 of the proof of Proposition~\ref{prop:ccc}; at \(F_1\), the stated entry is nondecreasing in the
posterior and \(\Delta_1>0\), so the prices differ. (E3): at \(F_0\), \(B_0(\pi)<c_-\) excludes the whole high
component, and \(B_0(\pi)>\bar c_0\) includes the whole low component. At \(F_1\), \(B_1(M)>c_+\) includes the whole high
component, and \(B_1(m)<B_1(\pi)\le B_0(\pi)<c_-\) excludes it at the other two posteriors. Every realized
comparison is strict, so both challenger tie rules agree. (E4): \eqref{eq:a2-terminal}. (E5): at \(F_0\), the correct gross
returns are \(\rho\Delta_0(1-\pi)\) and \(\rho\Delta_0\pi\), both below \(k\); at \(F_1\), the purchase earns
\(\Delta_1(1-M)-k>0\) and the sale \(\rho\Delta_1m-k>0\). (E6): truthful bids.

\emph{Step 4: outcomes.} A purchase has probability \(b+\lambda\) in class \(H\) and \(b\) in class \(L\), and only
a purchase admits the high component. So \(\Ent_1=\rho+(1-\rho)(b+\pi\lambda)\), class entry is
\(\bar e_H=\rho+(1-\rho)(b+\lambda)\), and ownership is \(\pi\bar e_HW_H(F_1)\) against \(\pi\rho W_H(F_0)\), as in
Step 4 of the proof of Proposition~\ref{prop:ccc}.
\end{proof}

The two-point law is the case \(\Phi^0=\delta_0\),
\(\Phi^1=\delta_{c_H}\).

\paragraph{The third economy and the example band.}
For \(F_2\), the conditions \(\bar c_0<B_2(m)\), \(B_2(M)<c_-\) and \(k<\rho\kappa\Delta_2\) give full orders by the same
bound, include the low component everywhere, and exclude the high component at every posterior, so entry is
\(\rho\). Suppose instead that the low component prepares everywhere, that \(B_r(\pi)<c_-\), and that the orders are
full. The high component then prepares only after a purchase, with probability \(\Phi^1(B_r(M))\), so
\[
\Ent(r)=\rho+(1-\rho)(b+\pi\lambda)\,\Phi^1\bigl(B_r(M)\bigr).
\]
This is the only place where atomlessness enters: for atomless \(\Phi^1\), \(\Phi^1\) is continuous, and \(B_r(M)\) falls
continuously in \(r\), so entry is continuous and nonincreasing in strength. In the example, take
\(\bar c_0=\reg{atom_L}\) and \(\Phi^1\) uniform on \([c_H-\varepsilon,c_H+\varepsilon]\) with \(\varepsilon=\reg{atom_eps}\).
The margins are \(B_1(m)-\bar c_0=\reg{margin_p4_atomless_floor_B}\), \(c_--B_0(\pi)=\reg{margin_p4_atomless_weak_B}\) and
\(B_1(M)-c_+=\reg{margin_p4_atomless_strong_B}\). Entry at \(r_1\) is \(\reg{p4_atomless_r1_entry}\) and ownership
is \(\reg{p4_atomless_r1_own}\). At \(r_2\), \(c_--B_2(M)=\reg{margin_p4_atomless_return_B}\): the top profit lies
inside the band, at distances \(\reg{margin_p4_atomless_top_lower_r2_B}\) and \(\reg{margin_p4_atomless_top_upper_r2_B}\)
from its ends, and the formula gives entry \(\reg{p4_atomless_r2_entry}\) and ownership \(\reg{p4_atomless_r2_own}\) by
direct evaluation of the uniform law. The return to \(\rho\) at \(r_2\) holds for every half-width below
\(c_H-B_2(M)=\reg{p4_atomless_max_width_floor_r2}\); at \(\varepsilon=\reg{atom_eps_small}\) its margin is
\(\reg{margin_p4_atomless_small_width_return}\).

\subsubsection*{Value classes and priors}

\paragraph{Value classes and priors.}\label{par:uniform-classes-scope}
For every pair of ranked finite-mean classes with a positive mean gap, a reserve \(p>0\), every interior prior and slot, and
\(0<\rho<1\), the uniform incumbent family \(F_r=U[0,r]\) admits a nonempty set, defined by strict inequalities
in strengths and costs, on which the hypotheses and ownership conclusion of Proposition~\ref{prop:ccc} hold.


\begin{proposition}[Value classes and priors]\label{prop:classes}
Fix classes that satisfy the class ranking~\eqref{eq:classes} with \(\E_H[\theta]>\E_L[\theta]\), a reserve \(p>0\), and
any \(\pi,\rho\in(0,1)\) and slot \((\lambda,\beta)\). For uniform incumbents, a nonempty open set of strengths and
costs \((r_0,r_1,r_2,c_H,k)\) satisfies (H1)--(H3) and the ownership condition of Proposition~\ref{prop:ccc}. On this
set every conclusion of Proposition~\ref{prop:ccc} holds.
\end{proposition}

\begin{proof}[Proof of Proposition~\ref{prop:classes}]
Let \(\Gamma(x)=\int_p^xD(y)\,dy\), and write \(\bar\Gamma=\E_H[\theta]-\E_L[\theta]>0\) for the total mean gap. Under
the class ranking~\eqref{eq:classes}, \(D\ge0\), so \(\Gamma\) is continuous and nondecreasing with limit \(\bar\Gamma\).
Let \(\theta_*=\inf\{x\ge p:\Gamma(x)>0\}\); finite means and \(\bar\Gamma>0\) give \(p\le \theta_*<\infty\).

\emph{Step 1: the uniform incumbent.} For \(R\sim U[0,r]\), \eqref{eq:spread} and integration by parts give
\[
\DT(r)=\int_p^r\Bigl(1-\frac yr\Bigr)D(y)\,dy=\frac1r\int_p^r\Gamma(y)\,dy ,
\]
since \(\Gamma(p)=0\) and the boundary term at \(r\) vanishes. Since \(\Gamma\) is continuous,
\[
\DT'(r)=\frac{r\Gamma(r)-\int_p^r\Gamma(y)\,dy}{r^2}\ge\frac{p\,\Gamma(r)}{r^2}>0\qquad(r>\theta_*),
\]
because \(\int_p^r\Gamma\le(r-p)\Gamma(r)\). Also \(\DT(r)\to0\) as \(r\downarrow \theta_*\), since \(\Gamma=0\) on
\([p,\theta_*]\). The split identity of Section~\ref{sec:profits} gives
\(g_H(r)-g_L(r)=\bar\Gamma-\DT(r)\ge\bar\Gamma-(r-p)\bar\Gamma/r=p\bar\Gamma/r>0\), and at \(r=\theta_*\), \(g_H-g_L=\bar\Gamma\). Dominated
convergence, with the integrable bound \(1-G_j\), gives continuity of \(g_j(r)\) and \(W_H(r)\) in \(r\), and
\(g_j(r)\to0\) as \(r\to\infty\). Since \(\theta\ge p\), \(W_H(r)=\E_H[\min\{\theta/r,1\}]\ge p/r>0\).

\emph{Step 2: the construction.} Let \(e_H=\rho+(1-\rho)(b+\lambda)>\rho\). At \(\theta_*\),
\(B_{\theta_*}(M)-B_{\theta_*}(\pi)=(M-\pi)\bar\Gamma>0\). By continuity choose \(r_1>\theta_*\) close enough to \(\theta_*\) that
\(B_{r_1}(M)>B_{\theta_*}(\pi)\) and \(e_HW_H(r_1)>\rho W_H(\theta_*)\). Then choose \(r_0\in(\theta_*,r_1)\) close enough to
\(\theta_*\) that \(\DT(r_0)<\rho\kappa\DT(r_1)\). Since \(B_r(\pi)\) and \(W_H(r)\) fall in \(r\),
\(B_{r_0}(\pi)\le B_{\theta_*}(\pi)<B_{r_1}(M)\) and \(e_HW_H(r_1)>\rho W_H(\theta_*)\ge\rho W_H(r_0)\). Choose \(c_H\) in
\((B_{r_0}(\pi),B_{r_1}(M))\) and \(k\) in \((\DT(r_0),\rho\kappa\DT(r_1))\). Since \(g_j(r)\to0\), choose
\(r_2>r_1\) with \(B_{r_2}(M)<c_H\). These choices satisfy (H1) to (H3) and the ownership condition.

\emph{Step 3: openness and the conclusions.} Each condition is a strict inequality between continuous functions of
\((r_0,r_1,r_2,c_H,k)\), so a neighborhood of the constructed point satisfies all of them. On this set,
The tuples of Proposition~\ref{prop:ccc} satisfy every equilibrium condition and are the unique outcomes. In
the point-mass case \(D=\1\{\ell\le y<h\}\), so \(\theta_*=\ell\).
\end{proof}

\begin{table}[tbp]
\centering\footnotesize\singlespacing
\caption{Value classes and priors for other class laws and priors.}\label{tab:classes}
\setlength{\tabcolsep}{4pt}
\begin{tabular}{@{}lllcccccc@{}}
\toprule
\multicolumn{9}{@{}l}{\emph{Panel A. Primitives}}\\
Economy & Low class & High class & \(\pi\) & \(c_H\) & \(k\) & \(r_0\) & \(r_1\) & \(r_2\)\\
\midrule
Moderate gap & \(\reg{p2_input_moderate_low_lo}\) & \(\reg{p2_input_moderate_high_lo}\) & \reg{p2_input_moderate_pi}
 & \reg{p2_input_moderate_cost} & \reg{p2_input_moderate_k} & \reg{p2_input_moderate_r0} & \reg{p2_input_moderate_r1} & \\
Narrow bands & \(U[\reg{p2_input_bands_low_lo},\reg{p2_input_bands_low_hi}]\) & \(U[\reg{p2_input_bands_high_lo},\reg{p2_input_bands_high_hi}]\)
 & \reg{p2_input_bands_pi} & \reg{p2_input_bands_cost} & \reg{p2_input_bands_k}
 & \reg{p2_input_bands_r0} & \reg{p2_input_bands_r1} & \reg{p2_input_bands_r2}\\
Low prior & \(\reg{p2_input_ex_a_low_lo}\) & \(\reg{p2_input_ex_a_high_lo}\) & \reg{p2_input_ex_a_pi}
 & \reg{p2_input_ex_a_cost} & \reg{p2_input_ex_a_k}
 & \reg{p2_input_ex_a_r0} & \reg{p2_input_ex_a_r1} & \reg{p2_input_ex_a_r2}\\
High prior & \(\reg{p2_input_ex_b_low_lo}\) & \(\reg{p2_input_ex_b_high_lo}\) & \reg{p2_input_ex_b_pi}
 & \reg{p2_input_ex_b_cost} & \reg{p2_input_ex_b_k}
 & \reg{p2_input_ex_b_r0} & \reg{p2_input_ex_b_r1} & \reg{p2_input_ex_b_r2}\\
Wide classes & \(U[\reg{p2_input_ex_c_low_lo},\reg{p2_input_ex_c_low_hi}]\) & \(U[\reg{p2_input_ex_c_high_lo},\reg{p2_input_ex_c_high_hi}]\)
 & \reg{p2_input_ex_c_pi} & \reg{p2_input_ex_c_cost} & \reg{p2_input_ex_c_k}
 & \reg{p2_input_ex_c_r0} & \reg{p2_input_ex_c_r1} & \reg{p2_input_ex_c_r2}\\
Same means, points & \(\reg{p2_input_ex_d_low_lo}\) & \(\reg{p2_input_ex_d_high_lo}\) & \reg{p2_input_ex_d_pi}
 & \reg{p2_input_ex_d_cost} & \reg{p2_input_ex_d_k}
 & \reg{p2_input_ex_d_r0} & \reg{p2_input_ex_d_r1} & \reg{p2_input_ex_d_r2}\\
\midrule
\multicolumn{9}{@{}l}{\emph{Panel B. Outcomes}}\\
Economy & Entry \(\Ent_0\) & Entry \(\Ent_1\) & \multicolumn{6}{l}{High-class ownership \(\Own_{H,1}\)}\\
\midrule
Moderate gap & \reg{p2_moderate_r0_entry} & \reg{p2_moderate_r1_entry} & \multicolumn{6}{l}{\reg{p2_moderate_r1_ownership}}\\
Narrow bands & \reg{p2_bands_r0_entry} & \reg{p2_bands_r1_entry} & \multicolumn{6}{l}{\reg{p2_bands_r1_ownership}}\\
Low prior & \reg{p2_ex_a_r0_entry} & \reg{p2_ex_a_r1_entry} & \multicolumn{6}{l}{\reg{p2_ex_a_r1_ownership}}\\
High prior & \reg{p2_ex_b_r0_entry} & \reg{p2_ex_b_r1_entry} & \multicolumn{6}{l}{\reg{p2_ex_b_r1_ownership}}\\
Wide classes & \reg{p2_ex_c_r0_entry} & \reg{p2_ex_c_r1_entry} & \multicolumn{6}{l}{\reg{p2_ex_c_r1_ownership}}\\
Same means, points & \reg{p2_ex_d_r0_entry} & \reg{p2_ex_d_r1_entry} & \multicolumn{6}{l}{\reg{p2_ex_d_r1_ownership}}\\
\bottomrule
\end{tabular}
\par\medskip
\begin{minipage}{0.95\linewidth}\footnotesize
\textit{Notes.} The slot with \(\beta=\reg{beta_B}\) and \(1+u=\eu\), the reserve \(\reg{prim_p}\) and \(\rho=\reg{prim_rho}\) in every row. A single
number is a point-mass class; \(U[\cdot,\cdot]\) is a uniform class. The moderate-gap economy has two strengths. At
the third strength every row with three strengths has entry \(\rho\). The moderate-gap row satisfies (H1)--(H2) of
Proposition~\ref{prop:ccc}; the other four reversal rows also satisfy (H3). The last row fails the recruitment condition at \(r_1\).
\end{minipage}
\end{table}

\paragraph{The economies of Table~\ref{tab:classes}.}
For a point value \(\theta\) and \(R\sim U[0,r]\), integration over the pieces cut by \(p\), \(\theta\) and \(r\) gives
\(t(\theta)=\theta-(\theta^2-p^2)/(2r)\), \(g(\theta)=(\theta^2-p^2)/(2r)\) and \(F(\theta)=\theta/r\) for \(p\le \theta\le r\), and \(t(\theta)=r/2+p^2/(2r)\),
\(g(\theta)=\theta-t(\theta)\) and \(F(\theta)=1\) for \(\theta\ge r\). A uniform class averages these piecewise polynomials over its support,
with a cut at \(r\); every piece has degree at most three. Exact integration gives (H1)--(H2) for the
moderate-gap row and (H1)--(H3) for the other four reversal rows, with the margins in Table~\ref{tab:margins}.
Proposition~\ref{prop:ccc} supplies their unique equilibrium outcomes. In the wide classes the class ranking
holds with overlapping supports: \(G_L-G_H\ge0\) everywhere, with equality only outside the common interior
of the supports.

The last row keeps the costs, the prior and the strengths of the wide classes and puts point masses at the same
means. Its upper margin at \(r_1\) is \(B_1(M)-c_H=\reg{margin_p2_ex_d_strong_top_B}<0\). The bounds of Steps 1 and 2
of the proof of Proposition~\ref{prop:ccc} still force zero orders at \(r_0\) and full orders at \(r_1\) and \(r_2\),
since their trade margins are positive (Table~\ref{tab:margins}). Then \(B_i(\mu)\le B_i(M)<c_H\) at every posterior
of the higher strengths, and the prior profit at \(r_0\) is below \(c_H\), so entry is \(\rho\) at all three strengths.
The comparison of the two rows is a convexity statement. For each \(R\), \(\theta\mapsto(\theta-\max\{p,R\})_+\) is convex, so a
mean-preserving increase in risk within a class weakly raises \(g_j\), and hence gross profit at every belief.

At the example's costs and strengths, the prior \(\reg{p2_input_ex_a_pi}\) gives entry
\(\reg{p2_fixed_prior_three_tenths_r1_entry}\) at all three strengths: the signed profit gap at \(r_1\) is
\(B_{r_1}(M(\reg{p2_input_ex_a_pi}))-c_H=\reg{p2_fixed_prior_three_tenths_r1_cost_sign_plus}\).
The prior \(\reg{p2_input_ex_b_pi}\) gives entry
\(\reg{p2_fixed_prior_seven_tenths_r0_entry}\) at \(r_0\), because the prior profit there exceeds the high cost by
\(\reg{p2_fixed_prior_seven_tenths_r0_cost_sign_plus}\), and entry \(\reg{p2_fixed_prior_seven_tenths_r1_entry}\) at
\(r_1\). In both cases the bound \(\DT(r_0)<k\) forces zero orders at \(r_0\), and the bound on every profile forces
full orders at \(r_1\) and \(r_2\).

\subsubsection*{More than two classes}

Let finitely many classes \(j\) have priors \(\pi_j>0\) and values above the reserve, and write \(t_{j,i}\) and
\(g_{j,i}\) for the class objects at \(F_i\). Every posterior vector \(\nu\) that a nonzero order can produce lies in the set
\[
\mathcal M_\pi=\Bigl\{\nu:\ \frac{\nu_j}{\pi_j}\le(1+u)\frac{\nu_l}{\pi_l}\ \text{for all } j,l\Bigr\}.
\]
The price reveals only the posterior mean of class proceeds, \(\nu\cdot t\). Partition the classes into a low group
\(\mathcal L\) and a high group \(\mathcal H\), with high-group prior mass \(\bar\pi\). Write \(\bar g_{\mathcal
H}\) and \(\bar g_{\mathcal L}\) for the prior-weighted group means of \(g_{j,1}\).

\begin{proposition}[Many classes under separation]\label{prop:manyclasses}
Let the incumbent strengthen from \(F_0\) to \(F_1\). Suppose
\[
\max_jt_{j,0}-\min_jt_{j,0}<k,\qquad \sum_j\pi_jg_{j,0}<c_H,
\]
\[
\min_{j\in\mathcal H}t_{j,1}-\max_{\nu\in\mathcal M_\pi}\nu\cdot t_1>\frac k\rho,\qquad
\min_{\nu\in\mathcal M_\pi}\nu\cdot t_1-\max_{j\in\mathcal L}t_{j,1}>\frac k\rho,
\]
and \(\bar g_{\mathcal H}>\bar g_{\mathcal L}\) with \(\bar g_{\mathcal L}+M(\bar\pi)(\bar g_{\mathcal H}-\bar
g_{\mathcal L})>c_H\). Then at \(F_0\) the unique orders are zero, and entry is \(\rho\). At \(F_1\) the unique orders
are a purchase in every high-group class and a sale in every low-group class, only a purchase recruits, and entry is
\(\rho+(1-\rho)(b+\bar\pi\lambda)>\rho\). These conclusions hold under either challenger rule.
\end{proposition}

\begin{proof}[Proof of Proposition~\ref{prop:manyclasses}]
The tuple lists every \(\sigma_j\), the posterior vector \(\mu(q)=(\mu_j(q))_j\) after every order,
the posterior vector at every price, the prices, \(\zeta\) and the bids. For \(q=\pm1\),
\(b\le f_j(q)\le b+\lambda\), so every ratio \(f_j(q)/f_l(q)\) lies in \([(1+u)^{-1},1+u]\).
Bayes' rule gives \(\mu_j(q)/\pi_j=f_j(q)/\sum_l\pi_lf_l(q)\), whence the posterior vectors at the two nonzero
orders belong to \(\mathcal M_\pi\) for every \(0<\beta<1/2\).

\emph{Payoffs.} Competitive prices give \(P(q)=t_0+e(q)\,[\mu(q)\cdot t-t_0]\), and
\(\E[V_T\mid j,q]=t_0+e(q)(t_j-t_0)\) as in \eqref{eq:a2-terminal}. Hence, against the fixed price and entry,
\[
U_j(q)=q\,e(q)\,[t_j-\mu(q)\cdot t]-k|q| ,
\]
with \(e(q)\in[\rho,1]\) because the zero-cost challenger prepares after every price. This step needs no price
sufficiency.

\emph{The weak economy.} Every residual satisfies \(|t_{j,0}-\mu\cdot t_0|\le\max_lt_{l,0}-\min_lt_{l,0}<k\), so
every nonzero order loses, and zero is the unique order of every class. Then \(f_j=f_l\) for all \(j,l\), every
posterior vector equals the prior vector \(\pi\), and gross profit is \(\sum_j\pi_jg_{j,0}<c_H\), so entry is \(\rho\).
The tuple has zero orders, the prior posterior at every order and price, the single price
\(t_{0,0}+\rho\,(\pi\cdot t_0-t_{0,0})\), \(\zeta\equiv0\), and truthful bids. Each condition holds by these
computations.

\emph{The strong economy, orders.} Under every candidate, a high-group class \(j\) earns from a purchase at least
\(\rho\,[\min_{l\in\mathcal H}t_{l,1}-\max_{\nu\in\mathcal M_\pi}\nu\cdot t_1]-k>0\), and from a sale at most \(-k\)
minus a nonnegative term, since \(t_{j,1}-\mu\cdot t_1>0\). A low-group class earns from a sale at least
\(\rho\,[\min_{\nu\in\mathcal M_\pi}\nu\cdot t_1-\max_{l\in\mathcal L}t_{l,1}]-k>0\), and from a purchase at most
\(-k\). Zero pays zero. So every equilibrium has purchases in \(\mathcal H\) and sales in \(\mathcal L\).

\emph{The strong economy, the tuple.} With these orders, the group posterior \(\mu_{\mathcal H}\) of \(\mathcal H\) follows the
two-class formulas with prior \(\bar\pi\): \(\mu_{\mathcal H}=(m(\bar\pi),\bar\pi,M(\bar\pi))\) after \((-1,0,1)\). Within each
group every class has the same order law, so the class proportions stay at the prior proportions:
\(\mu_j(\mu_{\mathcal H})=\mu_{\mathcal H}\pi_j/\bar\pi\) for \(j\in\mathcal H\) and \(\mu_j(\mu_{\mathcal H})=(1-\mu_{\mathcal H})\pi_j/(1-\bar\pi)\) for \(j\in\mathcal L\). Gross
profit at \(\mu_{\mathcal H}\) is \(\bar g_{\mathcal L}+\mu_{\mathcal H}(\bar g_{\mathcal H}-\bar g_{\mathcal L})\), increasing in \(\mu_{\mathcal H}\). At
\(\mu_{\mathcal H}=\bar\pi\) it is \(\sum_j\pi_jg_{j,1}\le\sum_j\pi_jg_{j,0}<c_H\) by Proposition~\ref{prop:returns}(a); at
\(\mu_{\mathcal H}=M(\bar\pi)\) it exceeds \(c_H\) by hypothesis. So only a purchase recruits, and entry is \((\rho,\rho,1)\). The
prices are \(P=t_{0,1}+e(\mu_{\mathcal H})[\bar t_{\mathcal L}-t_{0,1}+(\bar t_{\mathcal H}-\bar t_{\mathcal L})\mu_{\mathcal H}]\) with group means
\(\bar t\). Separation gives \(\bar t_{\mathcal H}>\max_{\nu}\nu\cdot t_1\ge\min_\nu\nu\cdot t_1>\bar t_{\mathcal L}\),
so with \(e\) nondecreasing in \(\mu_{\mathcal H}\) and \(e\ge\rho\), the three prices differ, each price reveals \(\mu_{\mathcal H}\) and the whole
vector, and \(\bar\mu(P)=\mu(\mu_{\mathcal H})\). Each condition then holds by direct computation: (E1) and (E2) by the formulas
above, (E3) by the strict profit comparisons, (E4) by the price formula, (E5) by the strict payoff bounds, (E6) by
truthful bids. A purchase has probability \(b+\bar\pi\lambda\), so entry is \(\rho+(1-\rho)(b+\bar\pi\lambda)>\rho\).
Every comparison is strict, so both challenger rules give this outcome.
\end{proof}

\paragraph{The four-class economy.}
The economy has the example slot, costs, reserve and first two strengths, and four point-mass classes with prior
\(\reg{four_prior}\) each. The low group has the values \(\reg{four_v1}\) and \(\reg{four_v2}\), and the high group has
the values \(\reg{four_v3}\) and \(\reg{four_v4}\).
For posterior vectors of the form \(\nu_j=\pi_j\iota_j/\sum_l\pi_l\iota_l\) with \(1\le \iota_j\le1+u\), the mean \(\nu\cdot t\) is
linear-fractional in each weight, so its extremes over \(\mathcal M_\pi\) occur at the vertices of the weight box.
Evaluation at every vertex gives the separation distance \(\varsigma=\reg{p2_groups_gamma}\), the trading margin
\(\rho\varsigma-k=\reg{p2_groups_trade}\), the weak range margin \(\reg{p2_groups_weak_range}\), the weak cost margin
\(\reg{p2_groups_weak_cost}\), the group profit gap \(\reg{p2_groups_profit_gap}\) and the upper profit margin
\(\reg{p2_groups_top}\).

\paragraph{A price collision at three classes.}
Without separation, the price can pool orders with different posterior vectors. Take the slot of the example,
\(R\sim U[0,\reg{a236_pool_r}]\), \(p=\reg{a236_pool_p}\), \(\rho=\reg{a236_pool_rho}\), \(c_H=\reg{a236_pool_cH}\), \(k=\reg{a236_pool_k}\), priors \((\reg{a236_pool_pi1},\reg{a236_pool_pi2},\reg{a236_pool_pi3})\), and the classes
\(G_1=\delta_p\), \(G_2=(1-\psi)\delta_p+\psi\delta_{\reg{a236_pool_h}}\), \(G_3=\delta_{\reg{a236_pool_h}}\) with \(\psi=(1+2b_0+\lambda)/2=\reg{p2_many_weight}\).
Since \(\beta>\reg{a236_pool_pi1}\), \(1-\psi=(1-\lambda)(4\beta-1)/2>0\), so the classes are strictly ordered. The auction gives \(t_0=p(1-p/r)\),
\[
t_1=p,\quad t_3=\frac r2+\frac{p^2}{2r},\quad t_2=p+\psi(t_3-p),\qquad
g_1=0,\quad g_3=h-t_3,\quad g_2=\psi g_3,
\]
where \(h=\reg{a236_pool_h}\) and \(r=\reg{a236_pool_r}\). The tuple is
\(\sigma_j=\delta_{j-2}\) (class 1 sells, class 2 abstains, class 3 buys), Bayes posteriors at each order, entry \(\rho\)
at every price, prices \(P(q)=t_0+\rho(\mu(q)\cdot t-t_0)\), the pooled belief at the common price of the orders \(0\)
and \(1\), and truthful bids. Every class profit is at most \(g_3<c_H\), and class 3 has positive probability at every
price, so the high cost never prepares and the zero cost prepares strictly. Normalize proceeds to \((0,\psi,1)\). The
conditional means at the orders \(0\) and \(1\) are
\[
\bar t(0)=\frac{b_0(\pi_2\psi+\pi_3)+\lambda\pi_2\psi}{b_0+\lambda\pi_2},\qquad
\bar t(1)=\frac{b(\pi_2\psi+\pi_3)+\lambda\pi_3}{b+\lambda\pi_3},
\]
and \(\bar t(0)=\bar t(1)\) exactly when \(\psi=(2b+3b_0+2\lambda)/[2(2b+b_0+\lambda)]\), which equals the stated \(\psi\) because
\(2b+b_0+\lambda=1\). The posterior vectors differ: \((\reg{p2_many_q1_mu0},\reg{p2_many_q1_mu1},\reg{p2_many_q1_mu2})\)
after the order \(0\) and \((\reg{p2_many_q2_mu0},\reg{p2_many_q2_mu1},\reg{p2_many_q2_mu2})\) after the order \(1\),
with the common price \(\reg{p2_many_q2_price}\). The sale has the distinct price \(\reg{p2_many_q0_price}\). Against
the fixed tuple, \(U_j(q)=q\rho[t_j-\mu(q)\cdot t]-k|q|\), and each class's chosen order beats its best alternative
by \(\reg{p2_many_choice0}\), \(\reg{p2_many_choice1}\) and \(\reg{p2_many_choice2}\). The collision is a knife-edge in
\(\psi\): a change of \(\reg{a236_pool_perturb}\) in either direction separates the two prices and keeps every choice strict.

\subsubsection*{Complementary private signals}

\begin{proposition}[Complementary private signals]\label{prop:signals}
Fix ranked finite-mean classes, a reserve \(p>0\), continuous laws \(F_1\le F_0\), prior one half, the
two-point preparation cost with \(0<\rho<1\), and a slot \(0<\lambda<1\), \(0<\beta<1/2\). Conditional on class
\(j\), the signals \(S_I,S_C\), the challenger value and the incumbent value are mutually independent, with
value laws \(G_j,F_i\). The cost and slot draws are independent of these variables. The binary signals have
accuracies \(1/2<\alpha_I,\alpha_C<1\):
\[
\Pr(S_I=+\mid H)=\Pr(S_I=-\mid L)=\alpha_I,\qquad
\Pr(S_C=+\mid H)=\Pr(S_C=-\mid L)=\alpha_C.
\]
The investor observes \(S_I\); the challenger observes \(S_C\) and the price. Define
\[
\begin{gathered}
M^I=(1-\alpha_I)+(2\alpha_I-1)M(1/2),\quad
m^I=(1-\alpha_I)+(2\alpha_I-1)m(1/2),\\
\mu^C_+(\mu)=\frac{\alpha_C\mu}{\alpha_C\mu+(1-\alpha_C)(1-\mu)},\quad
\mu^C_-(\mu)=\frac{(1-\alpha_C)\mu}{(1-\alpha_C)\mu+\alpha_C(1-\mu)},\\
\bar\Delta_P(F)=\DT(F)+(1-\rho)(2\alpha_C-1)[t_H(F)-t_0(F)].
\end{gathered}
\]
Suppose that
\[
B_{F_0}(\alpha_C)<c_H<B_{F_1}(\mu^C_+(M^I)),\qquad
(2\alpha_I-1)\bar\Delta_P(F_0)<k<(2\alpha_I-1)\rho\,m(1/2)\DT(F_1).
\]
Under all four tie rules, the unique orders at \(F_0\) are abstention after either investor signal, with entry
\(\rho\) and high-class ownership \(\rho W_H(F_0)/2\). At \(F_1\), the unique orders are a purchase after \(+\)
and a sale after \(-\); the public posteriors after \((-1,0,1)\) are \((m^I,1/2,M^I)\), the three prices differ,
and entry strictly exceeds \(\rho\).

At \(F_1\), write \(\zeta_\pm(\mu)\) for high-cost preparation following \(S_C=\pm\) at public belief \(\mu\).
For \(\mu^C_\pm(\mu)\) strictly above or below the preparation threshold it is respectively one or zero; at
equality the chosen challenger tie rule applies. Conditional entry is
\[
\begin{aligned}
e_H(\mu)&=\rho+(1-\rho)[\alpha_C\zeta_+(\mu)+(1-\alpha_C)\zeta_-(\mu)],\\
e_L(\mu)&=\rho+(1-\rho)[(1-\alpha_C)\zeta_+(\mu)+\alpha_C\zeta_-(\mu)].
\end{aligned}
\]
Its class average \(\bar e_H=\sum_x f_H(x)e_H(\mu(x))\) gives ownership \(\Own_{H,1}=\bar e_HW_H(F_1)/2\), where
\[
f_H=(b+(1-\alpha_I)\lambda,b_0,b+\alpha_I\lambda),\qquad
f_L=(b+\alpha_I\lambda,b_0,b+(1-\alpha_I)\lambda).
\]
If also \(B_{F_1}(\mu^C_-(M^I))<c_H\), exactly the pair \((X,S_C)=(1,+)\) recruits, and
\[
\bar e_H=\rho+(1-\rho)\alpha_C(b+\alpha_I\lambda),\qquad
\bar e_L=\rho+(1-\rho)(1-\alpha_C)[b+(1-\alpha_I)\lambda].
\]
In general high-class ownership rises exactly when \(\bar e_HW_H(F_1)>\rho W_H(F_0)\); in particular it rises if
\(W_H(F_1)=1\), or if decisive advantage holds at \(F_1\).
\end{proposition}

\begin{proof}
\emph{Payoff bounds.} The cost window and stochastic dominance imply
\(B_{F_1}(\mu^C_+(M^I))>B_{F_1}(\alpha_C)\), hence \(g_H(F_1)>g_L(F_1)\).
The profit gap at \(F_0\) is at least as large by Lemma~\ref{lem:twonumbers}, and the trading window gives
\(\DT(F_1)>0\). For a candidate, let
\(f_\pm(x)=(1-\lambda)\eta(x)+\lambda\sigma_\pm(x)\) and
\(\nu_I(x)=f_+(x)/[f_+(x)+f_-(x)]\). The two investor signals have equal prior probabilities, and conditional
independence gives
\[
\mu(x)=(1-\alpha_I)+(2\alpha_I-1)\nu_I(x).
\]
At either nonzero order Lemma~\ref{lem:beliefs}(b), applied to the investor signal, gives
\(m(1/2)\le\nu_I(x)\le M(1/2)\). At a price the challenger combines its signal with the public belief using
\(\mu^C_\pm\) from the statement. Write its preparation probabilities there as \(\zeta_\pm\).
The signal likelihood ratios are strictly ordered, so \(\zeta_+\ge\zeta_-\), under either challenger rule.
Consequently \(e_H\ge e_L\ge\rho\), and
\[
e_H-e_L=(1-\rho)(2\alpha_C-1)(\zeta_+-\zeta_-).
\]
With \(w_j=t_j-t_0\), conditional independence gives terminal payoff \(t_0+e_jw_j\) in class \(j\).
Competitive pricing is therefore \(P(x)=t_0+e_Lw_L+\mu(x)\Delta_P\), where
\begin{equation}\label{eq:a6-DP}
\Delta_P=e_Hw_H-e_Lw_L=e_L\DT+(e_H-e_L)w_H,\qquad
\rho\DT\le\Delta_P\le\bar\Delta_P.
\end{equation}
These bounds apply at every candidate price. If two orders share a price, the two signal-contingent preparation
probabilities agree, so their class entries agree. At \(F_1\), the positive lower bound in
\eqref{eq:a6-DP} then gives equal public posteriors. Bayes' rule at that price gives the common posterior.

Subtracting the price from the terminal payoff after each investor signal gives
\[
\begin{aligned}
U_+(1)&=(2\alpha_I-1)[1-\nu_I(1)]\Delta_P(1)-k,\\
U_-(-1)&=(2\alpha_I-1)\nu_I(-1)\Delta_P(-1)-k.
\end{aligned}
\]
The wrong-signed orders earn respectively
\(-(2\alpha_I-1)[1-\nu_I(-1)]\Delta_P(-1)-k\) and
\(-(2\alpha_I-1)\nu_I(1)\Delta_P(1)-k\), at most \(-k\). Zero pays zero.
Every correct gross return is at most \((2\alpha_I-1)\bar\Delta_P\) and at least
\((2\alpha_I-1)m(1/2)\rho\DT\).

\emph{The weaker law.} The upper trading bound at \(F_0\) is below \(k\), so abstention after either signal
is forced under both investor rules. All order and price posteriors are one half. The two private posteriors
are \(\alpha_C\) and \(1-\alpha_C\); both profits are below \(c_H\). Thus entry is \(\rho\), all orders share
price \(t_0+\rho(w_H+w_L)/2\), and truthful bids complete the tuple. The displayed payoff bounds give strict
optimality of abstention. Preparation is independent of both values given class, so high-class ownership is
\(\rho W_H(F_0)/2\).

\emph{The stronger law.} The lower trading bound at \(F_1\) exceeds \(k\), so the unique orders are a purchase
after \(S_I=+\) and a sale after \(S_I=-\). Then
\(f_+=(b,b_0,b+\lambda)\), \(f_-=(b+\lambda,b_0,b)\),
\(\nu_I=(m(1/2),1/2,M(1/2))\), and the public posteriors are \((m^I,1/2,M^I)\).
Choose \(\zeta_\pm\) by the prescribed preparation rule at these beliefs. Each is nondecreasing across the
three distinct beliefs, so both \(e_H\) and \(e_L\) are nondecreasing. If \(\mu_2>\mu_1\),
\[
\begin{aligned}
P(\mu_2)-P(\mu_1)
={}&(\mu_2-\mu_1)\Delta_P(\mu_1)\\
&+\mu_2[e_H(\mu_2)-e_H(\mu_1)]w_H\\
&+(1-\mu_2)[e_L(\mu_2)-e_L(\mu_1)]w_L>0.
\end{aligned}
\]
The first term is positive and the others nonnegative. The prices are therefore distinct, each reveals its
public posterior, and combining it with the private signal gives the stipulated preparation. The price formula,
strict optimal-order bounds and truthful bids establish the other equilibrium conditions.
At the purchase price \(B_{F_1}(\mu^C_+(M^I))>c_H\); the event \((X,S_C)=(1,+)\) has positive probability
in each class. Thus each class average of entry, and total entry, strictly exceeds \(\rho\).

\emph{Ownership and the recruiting pair.} Conditional on class, the signals, slot and cost select preparation
independently of \((\theta,R)\). Therefore
\(\Own_{H,1}=\bar e_HW_H(F_1)/2\), with the class-specific order laws in the statement.
If \(B_{F_1}(\mu^C_-(M^I))<c_H\), the unfavorable challenger signal never recruits.
After a sale or abstention, the public belief is at most one half, and
\(B_{F_1}(\mu^C_+(1/2))=B_{F_1}(\alpha_C)\le B_{F_0}(\alpha_C)<c_H\).
Thus exactly \((X,S_C)=(1,+)\) recruits. Its probabilities are
\(\alpha_C(b+\alpha_I\lambda)\) in class \(H\) and
\((1-\alpha_C)[b+(1-\alpha_I)\lambda]\) in class \(L\), giving the stated averages.
Comparing ownership at the two laws gives the exact condition. If \(W_H(F_1)=1\), dominance gives
\(W_H(F_0)=1\) and \(\bar e_H>\rho\) gives the strict rise. Decisive advantage implies this outbidding condition
by Lemma~\ref{lem:decisive}.
\end{proof}

\paragraph{A slot with a more accurate challenger signal.}
Take point classes \(h=\reg{prim_h}\), \(\ell=\reg{prim_ell}\), reserve \(p=\reg{prim_p}\), uniform incumbents on \([0,r]\),
\((r_0,r_1)=(\reg{sig_r0},\reg{sig_r1})\), and \((\rho,c_H,k)=(\reg{sig_rho},\reg{sig_cH},\reg{sig_k})\).
Let \(m=m(1/2)\), \(M=M(1/2)\), and set
\((\lambda,\beta)=(\reg{a236_sig_lambda},\reg{a236_sig_beta})\).
Then \((\alpha_I,\alpha_C)=(\reg{sig_aT},\reg{sig_d})\), \(M=\reg{a236_sig_M}\) and \(m=\reg{a236_sig_m}\). The five margins are
\[
\begin{aligned}
c_H-B_{r_0}(\alpha_C)&=\reg{a236_sig_alt_private}>0,\\
B_{r_1}(\mu^C_+(M^I))-c_H&=\reg{a236_sig_alt_joint}>0,\\
k-(2\alpha_I-1)\bar\Delta_P(r_0)&=\reg{a236_sig_alt_weak}>0,\\
(2\alpha_I-1)\rho m\DT(r_1)-k&=\reg{a236_sig_alt_strong}>0,\\
c_H-B_{r_1}(\mu^C_-(M^I))&=\reg{a236_sig_alt_other}>0.
\end{aligned}
\]
Thus the hypotheses allow \(\alpha_C>\alpha_I\).

\paragraph{The signal example.}
In the signal example at the slot \(\beta=\reg{beta_B}\), \(1+u=\eu\), with \((\alpha_I,\alpha_C)=(\reg{sig_aT},\reg{sig_d})\), the four conditions of the proposition and the one-pair
condition have the margins \(\reg{margin_p4_signal_s_private}\), \(\reg{margin_p4_signal_s_joint}\),
\(\reg{margin_p4_signal_s_weak}\), \(\reg{margin_p4_signal_s_strong}\) and \(\reg{margin_p4_signal_s_other_pair}\). The
formulas above give entry \(\reg{p4_s_signal_r0_entry}\) and \(\reg{p4_s_signal_r1_entry}\), and ownership
\(\reg{p4_s_signal_r0_own}\) and \(\reg{p4_s_signal_r1_own}\). Without the challenger's signal (\(\alpha_C=\reg{prim_pi}\)), the top
public profit stays below \(c_H\) at both strengths, with margins \(\reg{margin_p4_s_no_y_top_r0}\) and
\(\reg{margin_p4_s_no_y_top_r1}\), so entry is \(\rho\); the bounds \(\reg{margin_p4_s_no_y_weak}\) and
\(\reg{margin_p4_s_no_y_strong}\) force zero and full orders, and entry is \(\reg{p4_s_no_y_r0_entry}\) at both
strengths. When the investor observes the class and the challenger has no signal, the margins
\(\reg{margin_p4_s_class_weak}\), \(\reg{margin_p4_s_class_strong}\), \(\reg{margin_p4_s_class_prior_r0}\) and
\(\reg{margin_p4_s_class_top_r1}\) give zero orders, full orders and entry \(\reg{p4_s_class_r0_entry}\) at both
strengths. On the grid of accuracies, six pairs satisfy the conditions, all with \(\alpha_C>\alpha_I\); the smallest margin among
them is \(\reg{margin_p4_grid_72_75_weak}\).

\paragraph{The example.}
In the example the conditions cannot hold with \(\alpha_C\ge \alpha_I\). The strong bound needs \(2\alpha_I-1>k/[m\rho\DT(r_1)]=\reg{p4_b_signal_z_min}\). If \(\alpha_C\ge \alpha_I\), then \(2\alpha_C-1\ge2\alpha_I-1\), and the weak gross bound is
\[ (2\alpha_I-1)\bar\Delta_P(r_0)\ge\frac{k}{m\rho\DT(r_1)}\left[\DT(r_0)+(1-\rho)\frac{k w_H(r_0)}{m\rho\DT(r_1)}\right].\]
The right side exceeds \(k\) by \(\reg{margin_p4_b_d_ge_a_impossible}\).
With an accurate investor signal and a weak challenger signal, \((\alpha_I,\alpha_C)=(\reg{sig_comp_aT},\reg{sig_comp_d})\), the
five margins are \(\reg{margin_p4_signal_comp_b_private}\), \(\reg{margin_p4_signal_comp_b_joint}\),
\(\reg{margin_p4_signal_comp_b_weak}\), \(\reg{margin_p4_signal_comp_b_strong}\) and
\(\reg{margin_p4_signal_comp_b_other_pair}\). Entry rises from \(\reg{p4_comp_b_r0_entry}\) to
\(\reg{p4_comp_b_r1_entry}\). Without the challenger's signal, the top public profit at \(r_1\) stays below \(c_H\),
with margin \(\reg{margin_p4_comp_b_no_y_top_r1}\), the bounds \(\reg{margin_p4_comp_b_no_y_weak}\) and
\(\reg{margin_p4_comp_b_no_y_strong}\) fix the orders, and entry is \(\reg{p4_comp_b_no_y_r1_entry}\) at both strengths.
Both sources are therefore needed for the rise. Every comparison is strict, so all four pairs of tie rules give the
same outcomes. At a higher investor accuracy, the public posterior alone can recruit, and a pair can satisfy the
theorem without this complementarity; the margin \(c_H-B_{r_1}(M^I)\) is the condition that separates the two.

\paragraph{The example with the signal example's accuracies.}
In the example with \((\alpha_I,\alpha_C)=(\reg{sig_aT},\reg{sig_d})\), three of the four conditions fail, and each equilibrium has a mixed order: both signal types mix at \(r_0\), and the
type with \(S_I=-\) mixes at \(r_1\).
After wrong-signed orders are excluded, write \(s_+=\lambda\sigma_+(1)\) and \(s_-=\lambda\sigma_-(-1)\); the purchase
return depends only on \(s_+\), and the sale return only on \(s_-\), with \(1-\nu_I(1)=b/(2b+s_+)\) and
\(\nu_I(-1)=b/(2b+s_-)\).
\begin{itemize}[leftmargin=2em]
\item \emph{At \(r_0\).} Under every profile \(B_{r_0}(\mu^C_+(m^I))>c_H\) and \(B_{r_0}(\mu^C_-(M^I))<c_H\), with margins
\(\reg{margin_p4_b_weak_all_yplus}\) and \(\reg{margin_p4_b_weak_no_yminus}\). So \(\zeta_+=1\) and \(\zeta_-=0\) at every price,
\(e_H=\reg{a236_sig_eH}\), \(e_L=\reg{a236_sig_eL}\), and \(\Delta_P=\reg{a236_sig_DP0}\) at every price. Each correct gross return is
\((2\alpha_I-1)\cdot\tfrac{b}{2b+s_\pm}\cdot\reg{a236_sig_DP0}\), strictly decreasing in its rate, above \(k\) at rate zero (margin
\(\reg{margin_p4_b_weak_zero_return}\)) and below \(k\) at rate \(\lambda\) (margin \(\reg{margin_p4_b_weak_full_loss}\)).
The unique zero-profit rates are \(s_+=s_-=\reg{a236_sig_rate0}b<\lambda\) (margin \(u-\reg{a236_sig_rate0}=\reg{margin_p4_b_weak_rate_room}\)). Entry is
\(\reg{p4_b_signal_r0_entry}\) and ownership \(\reg{p4_b_signal_r0_own}\).
\item \emph{At \(r_1\).} A challenger with \(S_C=-\) never recruits (margin \(\reg{margin_p4_b_strong_no_yminus}\)). A
purchase has public posterior at least \(1/2\), and \(B_{r_1}(\alpha_C)>c_H\) (margin \(\reg{margin_p4_b_strong_prior_yplus}\)),
so \(S_C=+\) recruits after a purchase, \(\Delta_P=\reg{a236_sig_DP1}\) there, and the purchase earns at least
\((2\alpha_I-1)m\reg{a236_sig_DP1}-k=\reg{margin_p4_b_strong_buy}>0\): \(s_+=\lambda\). After a sale, \(S_C=+\) recruits exactly when
\(\mu(-1)\ge\reg{a236_sig_public_cutoff}\), that is \(s_-\le\reg{a236_sig_sale_cutoff}b\). On that range the sale earns at least the return at the public belief that makes a challenger with \(S_C=+\) indifferent,
\(\reg{a236_sig_sale_floor}\), which exceeds \(k\) by \(\reg{margin_p4_b_strong_threshold_floor_return}\); so no equilibrium has
\(s_-\le\reg{a236_sig_sale_cutoff}b\), and no tie mixture exists at the threshold. Above it, \(e_H=e_L=\rho\) at the sale price.
The net sale payoff \(U_-(-1)=(2\alpha_I-1)\rho\DT(r_1)b/(2b+s_-)-k\) strictly decreases with \(s_-\).
Its unique zero-profit rate is \(s_-=\reg{a236_sig_rate1}b<\lambda\) (margin
\(\reg{margin_p4_b_strong_rate_room}\)), with \(\mu(-1)=\reg{a236_sig_sale_mu}\), below \(\reg{a236_sig_public_cutoff}\) by \(\reg{margin_p4_b_strong_sale_public_gap}\).
After the zero order \(S_C=+\) recruits (margin \(\reg{margin_p4_b_strong_zero_yplus}\)). Only a sale blocks recruitment
after \(S_C=+\); a sale has probability \(b+(1-\alpha_I)s_-\) in class \(H\) and \(b+\alpha_I s_-\) in class \(L\). Thus \(\bar e_H=\rho+(1-\rho)\alpha_C[1-b-(1-\alpha_I)s_-]\) and
\(\bar e_L=\rho+(1-\rho)(1-\alpha_C)[1-b-\alpha_I s_-]\). Hence entry is
\((\bar e_H+\bar e_L)/2=\reg{p4_b_signal_r1_entry}\) and ownership \(\bar e_H/2=\reg{p4_b_signal_r1_own}\).
\end{itemize}
At \(r_0\), the investor mixes between trade and abstention after either signal, and it trades with probability
\(\reg{p4_b_signal_r0_buy_rate}\). At \(r_1\), it buys with probability one after a favorable signal, and after an
unfavorable signal it sells with probability \(\reg{p4_b_signal_r1_sell_rate}\).
Each mixed tuple is unique when an indifferent investor may mix, with either challenger rule, by the strict monotonicity of each return and
the endpoint signs. When an indifferent investor abstains, neither strength has an equilibrium, because the unique qualifying rates are
interior. Without the challenger's signal, entry is \(\rho\) under every profile at both strengths, zero orders are
forced at \(r_0\), and at \(r_1\) the unique rates are \(s_+=s_-=\reg{a236_sig_rate1}b\); this control also requires that an indifferent investor may mix, since the gross
return at zero, \(\reg{p4_b_no_y_r1_zero_gross}\), exceeds \(k\) and the return at full orders,
\(\reg{p4_b_no_y_r1_full_gross}\), falls short of it.
